% Generated draft: requires mathematical and source review.
\documentclass[11pt]{article}
\usepackage[margin=25mm]{geometry}
\usepackage{fontspec}
\setmainfont{Latin Modern Roman}
\usepackage{amsmath,amssymb,mathtools,unicode-math}
\setmathfont{Latin Modern Math}
\usepackage{enumitem,needspace,xurl,hyperref}
\setlength{\parindent}{0pt}
\setlength{\parskip}{5pt}
\setlength{\emergencystretch}{3em}
\newcommand{\N}{\mathbb{N}}
\newcommand{\Z}{\mathbb{Z}}
\newcommand{\Q}{\mathbb{Q}}
\newcommand{\R}{\mathbb{R}}
\newcommand{\C}{\mathbb{C}}
\newcommand{\F}{\mathbb{F}}
\newcommand{\A}{\mathbb{A}}
\newcommand{\PP}{\mathbb{P}}
\newcommand{\E}{\mathbb{E}}
\newcommand{\eps}{\varepsilon}
\newcommand{\dd}{\,\mathrm{d}}
\begin{document}
\section*{UnsolvedMath ID: 6700045}
\textbf{Source problem code:} catalogue entry
\subsection{Definitions and mathematical statement}
A Riemannian manifold is complete when its distance metric is complete. The source asks: when is ``such an $X_0$'' area extremal in the category of complete manifolds? The preceding description of $X_0$ is absent, so its metric, curvature assumptions, and topology are not recoverable from this item.

For context, in a same-manifold metric formulation one compares a reference metric $g_0$ with complete metrics $g$ satisfying
\[
|v\wedge w|_g\geq|v\wedge w|_{g_0},\qquad
\operatorname{Sc}_g\geq\operatorname{Sc}_{g_0},
\]
and asks whether scalar-curvature equality is forced. Map-based versions can instead compare different manifolds through area-nonincreasing maps and additional topological data. The source does not specify which version of area extremality is intended, so neither is imposed as its exact definition.

Determine conditions on the intended $X_0$ that ensure the stated extremality against the intended complete competitors. Completeness does not by itself specify behavior at infinity, comparison-map properness, or asymptotic agreement. These hypotheses and the missing antecedent must be supplied before one can formulate a precise rigidity/extremality theorem. The source does not explicitly exclude compact manifolds from this complete category.
\subsection{Short English statement}
Determine when the reference manifold from the missing context remains area extremal among complete competitors. Its hypotheses and the exact area-comparison notion are not included in the extract.
\subsection{Sources}
\begin{itemize}
\item Misha Gromov, \emph{101 Questions, Problems and Conjectures around Scalar Curvature} (2017), printed tag [?48], occurrence 47, PDF page 54, \url{https://www.ihes.fr/~gromov/wp-content/uploads/2018/08/101-problemsOct1-2017.pdf}; accessed 2026-07-29, PDF-text-normalized extraction. The provenance describes the public author PDF as incomplete.
\item The supplied summary mentions complete-manifold rigidity, Brendle--Marques--Neves-type comparisons, and rigidity at infinity, without precise theorem citations or hypotheses.
\item Local registry \texttt{data/unsolvedmath/problems.json}, permanent UnsolvedMath ID $6700045$, source code catalogue entry, \emph{Scalar Curvature Question [?48]}. Dataset attribution: Ulam AI and the UnsolvedMath Contributors.
\end{itemize}
\end{document}
